\documentclass[10pt,a4paper]{article}

\usepackage[T1]{fontenc}
\usepackage[utf8]{inputenc}
\usepackage{tgtermes}
\usepackage{tgheros}
\usepackage{microtype}
\usepackage[a4paper,margin=18mm,headheight=19pt]{geometry}
\usepackage{xcolor}
\usepackage{graphicx}
\usepackage{booktabs}
\usepackage{tabularx}
\usepackage{array}
\usepackage{enumitem}
\usepackage[most]{tcolorbox}
\usepackage{fancyhdr}
\usepackage{lastpage}
\usepackage{titlesec}
\usepackage[hidelinks]{hyperref}
\usepackage{xurl}
\usepackage{caption}
\usepackage{placeins}
\usepackage{textcomp}
\usepackage{amsmath}
\usepackage{longtable}

\definecolor{Navy}{HTML}{173753}
\definecolor{Teal}{HTML}{008F8C}
\definecolor{PaleTeal}{HTML}{E5F4F3}
\definecolor{Orange}{HTML}{F08A24}
\definecolor{PaleOrange}{HTML}{FFF1E2}
\definecolor{Red}{HTML}{BB3E35}
\definecolor{PaleRed}{HTML}{F9E9E7}
\definecolor{Slate}{HTML}{5B7083}
\definecolor{Light}{HTML}{F3F6F8}
\definecolor{Line}{HTML}{D7E0E7}
\definecolor{Green}{HTML}{2E8B57}
\definecolor{PaleGreen}{HTML}{EAF5EE}

\renewcommand{\familydefault}{\rmdefault}
\newcommand{\sans}{\fontfamily{qhv}\selectfont}
\newcommand{\reportdate}{31 August 2026}
\newcommand{\reportversion}{Version 8}
\newcommand{\fom}{\mathrm{FOM}}
\newcommand{\secpd}{\mathrm{SEC}}
\newcommand{\evidencepath}[1]{\path{#1}}
\newcommand{\reportfigure}[3][0.94]{%
  \begin{figure}[htbp]
    \centering
    \includegraphics[width=#1\textwidth]{#2}
    \caption{#3}
  \end{figure}%
}

\hypersetup{
  pdftitle={HeiQ XMembrane Current Team Result and Decision Report -- Version 8},
  pdfauthor={HeiQ XMembrane project team},
  pdfsubject={Integrated model, material and coupon-cell decision basis as of 31 August 2026},
  pdfkeywords={XMembrane, electro-osmosis, ERV, O30, O50, membrane, material, coupon cell, FOM, decision report}
}

\pagestyle{fancy}
\fancyhf{}
\fancyhead[L]{\sans\footnotesize\color{Slate}HEIQ MATERIALS AG · XMEMBRANE}
\fancyhead[R]{\sans\footnotesize\color{Slate}TEAM RESULT REPORT · V8}
\fancyfoot[L]{\sans\scriptsize\color{Slate}Internal modelling and experiment-planning evidence · no product validation}
\fancyfoot[R]{\sans\scriptsize\color{Slate}Page \thepage\ of \pageref*{LastPage}}
\renewcommand{\headrulewidth}{0.45pt}
\renewcommand{\headrule}{\hbox to\headwidth{\color{Navy}\leaders\hrule height \headrulewidth\hfill}}
\renewcommand{\footrulewidth}{0pt}

\titleformat{\section}{\sans\bfseries\color{Navy}\fontsize{19}{22}\selectfont}{\thesection}{0.65em}{}
\titleformat{\subsection}{\sans\bfseries\color{Navy}\fontsize{13}{16}\selectfont}{\thesubsection}{0.55em}{}
\titleformat{\subsubsection}{\sans\bfseries\color{Teal}\fontsize{11}{13}\selectfont}{\thesubsubsection}{0.45em}{}
\titlespacing*{\section}{0pt}{7mm}{3mm}
\titlespacing*{\subsection}{0pt}{5mm}{2mm}
\titlespacing*{\subsubsection}{0pt}{4mm}{1.5mm}

\setlength{\parindent}{0pt}
\setlength{\parskip}{1.6mm}
\setlist[itemize]{leftmargin=5mm,itemsep=1mm,topsep=1mm}
\setlist[enumerate]{leftmargin=6mm,itemsep=1.2mm,topsep=1mm}
\captionsetup{font=small,labelfont={bf,color=Navy}}
\renewcommand{\arraystretch}{1.18}
\newcolumntype{Y}{>{\raggedright\arraybackslash}X}
\newcolumntype{R}{>{\raggedleft\arraybackslash}X}
\newcolumntype{C}{>{\centering\arraybackslash}X}

\newtcolorbox{mainbox}{
  enhanced,breakable,colback=PaleTeal,colframe=Teal,boxrule=0.7pt,
  arc=2mm,left=4mm,right=4mm,top=3mm,bottom=3mm
}
\newtcolorbox{warningbox}{
  enhanced,breakable,colback=PaleOrange,colframe=Orange,boxrule=0.7pt,
  arc=2mm,left=4mm,right=4mm,top=3mm,bottom=3mm
}
\newtcolorbox{stopbox}{
  enhanced,breakable,colback=PaleRed,colframe=Red,boxrule=0.7pt,
  arc=2mm,left=4mm,right=4mm,top=3mm,bottom=3mm
}
\newtcolorbox{gobox}{
  enhanced,breakable,colback=PaleGreen,colframe=Green,boxrule=0.7pt,
  arc=2mm,left=4mm,right=4mm,top=3mm,bottom=3mm
}
\newtcolorbox{plainbox}{
  enhanced,breakable,colback=Light,colframe=Line,boxrule=0.5pt,
  arc=2mm,left=4mm,right=4mm,top=3mm,bottom=3mm
}

\newcommand{\eyebrow}[1]{%
  {\sans\bfseries\color{Teal}\MakeUppercase{#1}}\par\vspace{1.5mm}%
}
\newcommand{\metric}[3]{%
  \begin{minipage}[t]{0.30\textwidth}
    {\sans\bfseries\fontsize{18}{20}\selectfont\color{Navy}#1}\par
    {\sans\bfseries\small #2}\par
    {\scriptsize\color{Slate}#3}
  \end{minipage}%
}

\begin{document}

\thispagestyle{empty}
\vspace*{7mm}
\begin{tcolorbox}[
  enhanced,colback=Navy,colframe=Navy,boxrule=0pt,arc=0mm,
  left=12mm,right=12mm,top=12mm,bottom=12mm
]
{\sans\bfseries\large\color{white}HEIQ XMEMBRANE PROJECT TEAM}\par
\vspace{8mm}
{\sans\bfseries\fontsize{28}{32}\selectfont\color{white}
Current result and decision\\[-1mm]
report}\par
\vspace{2mm}
{\sans\fontsize{18}{21}\selectfont\color{white}
Integrated system, membrane, material and coupon-cell evidence}\par
\vspace{7mm}
{\sans\large\color{white}\reportdate\quad · \quad\reportversion}
\end{tcolorbox}

\vspace{10mm}
{\sans\bfseries\Large\color{Navy}Decision basis for the project team}\par
\vspace{4mm}

This report consolidates the current model closure, the post-ERV full-duty
and percentage studies, common-voltage results, system and figure-of-merit
ledgers, membrane/material screens, and the centimetre-scale vapour-cell
design. It supersedes earlier team-facing interpretations for current
decisions, while retaining their machine-readable evidence as provenance.

\begin{mainbox}
\textbf{Current milestone.} The broad modelling phase has reached its stop
criterion. The next phase is material and subsystem feasibility: acquire a
same-sample membrane-property bundle, build the controlled-flow coupon cell,
and close the receiver/thermal/hardware ledger. No material, module,
technology or net full-system saving has been validated.
\end{mainbox}

\vfill
\begin{tabularx}{\textwidth}{@{}lY@{}}
\textbf{Primary equal-service reference} & Passive ERV plus remaining AC:
543.741 kW at the declared O80 outdoor-air duty \\
\textbf{Leading first experiment} & Post-ERV EO toward O30 in a controlled-flow
25-cm\textsuperscript{2} coupon cell \\
\textbf{Architecture status} & O30 is sampled below 2 V on one common bus;
full remaining post-ERV duty is not demonstrated on one common bus \\
\textbf{Release status} & Model-phase delivery is internally checkable;
immutable release snapshot remains open because the working tree is not frozen \\
\textbf{Claim status} & Conditional model and experiment-planning evidence only \\
\end{tabularx}

\clearpage
\pdfbookmark[0]{Abstract}{abstract}
\section*{Abstract}
\addcontentsline{toc}{section}{Abstract}

The present evidence supports one clear phase decision: stop broad parameter
expansion and begin controlled material and subsystem testing. At the
150,000-cfm O80 duty, the equal-service electrical reference is
543.741~kW for passive ERV plus the remaining AC load, compared with
1,166.4~kW for AC alone. A post-ERV, common-voltage O30 coordinate gives a
known subtotal of 381.327~kW and therefore a conditional 162.414-kW window
relative to passive ERV+AC. This is not a realised saving because receiver/sink
energy, EO heat rejection, electrodes, manifolds, leakage, durability and
measured material performance remain open. O50 requires no post-ERV EO duty
on the declared absolute-humidity basis; O30 closes 72.700\% of the remaining
post-ERV moisture gap and is the strongest first experimental target.

The full-duty and percentage studies show that the historical 10\%
coordinate actually closes 22.193\% of the remaining post-ERV duty. Uniform-
current sensitivity rows can close full duty below 2~V, but the independent
one-common-bus replay has no confirmed full-duty point on the sampled
0.5--2.0-V grid; 21 refined rows are physical no-solutions and 28 remain
numerical/local unresolved. This is not a global no-go proof. At O30,
sampled common-bus feasibility survives from $D=10^{-14}$ to
$10^{-9}$~m\textsuperscript{2}/s, but adverse diffusion erodes the energy
window and reverses it at the highest sampled diffusivity.

The measured-data requirement is a coupled, same-sample bundle: sorption and
hysteresis, chemical water mobility/permeance, electro-osmotic drag, wet
conductivity and compressed ASR, thickness/swelling/mechanics, coupled net
water flux, and receiver/electrode compatibility. Nafion is a useful numerical
reference but does not match the required drag-property vector and is not a
validated candidate. The proposed first fixture uses a 50$\times$50-mm active
aperture, an approximately 80$\times$80-mm coupon, 110$\times$110-mm outer
plates, interchangeable 1/2/4-mm gaps and controlled flow. Polarity reversal,
two-sided water balance, pressure control, blanks and gap/flow repeats are
required to separate EO transport from static humidity/temperature mismatch,
films, edge leakage and wall inventory.

\vspace{4mm}
\begin{tabularx}{\textwidth}{@{}CCC@{}}
\metric{162.4 kW}{Conditional O30 window}{Known common-voltage subtotal before open losses} &
\metric{72.70\%}{Remaining duty at O30}{O50 requires no post-ERV EO duty} &
\metric{25 cm\textsuperscript{2}}{First active aperture}{Controlled flow, ABBA and two-sided balance}
\end{tabularx}

\clearpage
\tableofcontents
\clearpage

\section{Executive conclusion and decision map}
\eyebrow{What is decided now}

\begin{gobox}
\textbf{GO: material and coupon work.} Begin a same-sample material search and
the controlled-flow vapour-cell experiment. Use O30 after the passive ERV as
the first system-relevant transport target. Continue modelling only to replay
measured candidate data, pre-register experiments, or close a named subsystem.
\end{gobox}

\begin{warningbox}
\textbf{OPEN: architecture and receiver.} Do not yet choose one-bus full duty,
voltage zoning, staged EO, direct EO inside the ERV, or a liquid/solid sink.
Those decisions require measured film behaviour, receiver work, heat
rejection, pressure drop, leakage, electrodes and durability.
\end{warningbox}

\begin{stopbox}
\textbf{NO-GO: validation language.} Do not state that 162.4~kW or any larger
model window has been saved, that full duty is impossible, that Nafion or an
inverse-design vector is a candidate membrane, or that the technology is
validated. The working tree must also be frozen before an immutable release
claim is made.
\end{stopbox}

\begin{center}\small
\begin{tabularx}{\textwidth}{@{}p{47mm}p{23mm}Y@{}}
\toprule
\textbf{Decision or question} & \textbf{Status} & \textbf{Current basis} \\
\midrule
Equal-service energy comparator & \textcolor{Green}{\textbf{CLOSED}} &
Use 543.741~kW passive ERV+AC as the primary reference \\
O30 common-voltage mechanism below 2 V & \textcolor{Green}{\textbf{SUPPORTED}} &
Confirmed sampled points, including 1.00~V, with local current sharing \\
Full post-ERV duty on one bus below 2 V & \textcolor{Orange}{\textbf{OPEN}} &
No confirmed sampled point; physical and unresolved rows coexist \\
Broad model phase & \textcolor{Green}{\textbf{STOP}} &
Core balances, signs, numerical replays and decision coordinates are sufficient for handoff \\
Material candidate ranking & \textcolor{Orange}{\textbf{NOT READY}} &
Zero of nine required same-sample data rows is complete \\
Receiver and thermal contract & \textcolor{Orange}{\textbf{OPEN}} &
Receiver/sink, disposal/regeneration and EO heat rejection are not closed \\
Practical module limits & \textcolor{Orange}{\textbf{OPEN}} &
Team must set area, pressure-drop, footprint, cost and lifetime limits \\
Net system or product claim & \textcolor{Red}{\textbf{NOT ALLOWED}} &
No measured full ledger or validated material/module \\
Material search and coupon cell & \textcolor{Green}{\textbf{GO}} &
Use the joint-property schema and controlled-flow boundary protocol \\
\bottomrule
\end{tabularx}
\end{center}

\subsection{Current integrity state}

The current final-delivery checker classifies the model-phase package as
deliverable and the material search as ready to start. It keeps the immutable
release snapshot open because the repository is dirty, and it keeps candidate
ranking, full-duty single-bus maturity, receiver closure and full-system claims
open. Audit PASS means declared balances, identities, hashes and schema checks
passed at that snapshot; it is not a technology pass.

\section{System architecture and equal-service comparison}
\eyebrow{ERV precedes EO in the leading post-ERV architecture}

The leading near-term chain is
\[
  \text{outdoor air}
  \longrightarrow \text{passive ERV}
  \longrightarrow \text{EO polishing}
  \longrightarrow \text{remaining AC sensible/latent service}.
\]
EO therefore acts on the finite moisture duty that remains after the ERV.
This is distinct from direct EO integrated into the ERV core, and from an
EO-first architecture without a passive latent exchanger.

\begin{center}\small
\begin{tabularx}{\textwidth}{@{}Y r r r Y@{}}
\toprule
\textbf{System} & \textbf{AC} & \textbf{Pre-dryer} & \textbf{Total} &
\textbf{Evidence boundary} \\
& [kW] & [kW] & [kW] & \\
\midrule
AC only & 1,166.400 & 0 & 1,166.400 &
Calibrated reference anchor \\
Passive ERV + AC & 509.381 & 34.359 & 543.741 &
Modelled passive equal-service reference \\
Passive ERV + post-ERV EO to O30 + AC & 316.910 & 64.417 & 381.327 &
\textbf{Conditional incomplete ledger}; receiver/sink and several hardware terms omitted \\
Direct active EO in ERV, full humidity duty + AC & 44.438 & 63.419 & 107.856 &
\textbf{Conditional architecture screen}; unmeasured inverse-design vector and incomplete coupling \\
\bottomrule
\end{tabularx}
\end{center}

\reportfigure[0.96]{../../outputs/fom_assessment_v1/figures/system_comparison.pdf}{
Current combined system comparison on one equal outdoor-to-supply water duty.
The AC-only and passive-ERV rows are references. Both EO rows are conditional
ledgers: they include the declared known electrical terms but exclude open
receiver/sink, thermal and hardware penalties.}
\FloatBarrier

Relative to passive ERV+AC, the current common-voltage O30 subtotal is lower
by 162.414~kW, or 29.87\%. Relative to AC alone it is lower by
785.073~kW. These differences are \emph{available conditional budgets}, not
measured savings. Every omitted positive load consumes the budget:
\[
P_{\mathrm{EO,all-in}}
=P_{\mathrm{grid}}+P_{\mathrm{module\,fan}}
+P_{\mathrm{receiver/sink}}+P_{\mathrm{heat\,rejection}}
+P_{\mathrm{BOP}}.
\]

\begin{plainbox}
\textbf{Why EO follows ERV here.} The passive ERV removes 3,137.65~kg/h from
the outdoor-to-supply water duty before EO is invoked. This reduces the EO
duty and makes O50 a zero-EO case. The direct active-ERV architecture remains
a model-level research alternative, but it has a different boundary,
different membrane preference and less subsystem closure.
\end{plainbox}

\section{O50, O30 and the meaning of percentage}
\eyebrow{Use finite remaining duty, not ambiguous relative humidity labels}

\subsection{O50}

At the frozen O80 outdoor state and passive-ERV reference, the ERV outlet is
already drier than the O50 absolute-humidity coordinate. O50 therefore
requires \textbf{no post-ERV EO duty}. The O50 bar is not evidence for an EO
saving; it is the passive-ERV reference state.

\subsection{O30}

O30 removes 1,228.27~kg/h after the passive ERV and closes 72.7003\% of the
finite moisture gap between the ERV outlet and the final supply humidity
ratio. It is the clearest first target because it:
\begin{itemize}
  \item creates a large but finite and interpretable water duty;
  \item has confirmed sampled one-common-bus solutions below 2~V;
  \item retains a sizeable conditional receiver/sink budget;
  \item exposes the decisive diffusivity, film and material interactions
        without requiring the full-duty pinch to close.
\end{itemize}

\subsection{The historical 10\% point}

The previously named 10\% reduction removes 374.94~kg/h. Relative to the
finite post-ERV-to-supply duty, it is \textbf{22.1925\% completion}. Both
coordinates are retained for traceability; the historical result must not be
silently relabelled.

\begin{center}\small
\begin{tabularx}{\textwidth}{@{}Y r r r r@{}}
\toprule
\textbf{Coordinate} & \textbf{EO water} & \textbf{Remaining-duty completion} &
\textbf{Fixed-trajectory area} & \textbf{Maximum voltage} \\
& [kg/h] & [\%] & [m\textsuperscript{2}] & [V] \\
\midrule
Historical 10\% & 374.94 & 22.1925 & 1,858.5 & 0.503 \\
O30 & 1,228.27 & 72.7003 & 6,089.1 & 1.314 \\
Full remaining duty & 1,689.50 & 100.0000 & 8,376.2 & 1.895 \\
\bottomrule
\end{tabularx}
\end{center}

The area/voltage values above are the fixed Terminal-1 uniform-current
sizing trajectory. They do not establish a common-bus module. The independent
common-voltage results are reported separately in Section~\ref{sec:commonbus}.

\section{Increasing EO share after the ERV}
\eyebrow{The energy trend is conditional and the endpoint changes control regime}

\reportfigure[0.96]{../../outputs/post_erv_eo_percentage_sweep_v1/percentage_energy_saving.png}{
Conditional energy effect as EO closes an increasing fraction of the
remaining post-ERV duty. The exact 100\% endpoint is a different coil-control
regime and must not be connected continuously to the one-sided near-endpoint
curve.}

Below the exact full-duty endpoint, the screening coil remains on its
saturated-coil branch. With zero EO process-heat leakage, the conservative
energy window grows approximately linearly with treated water when all-in EO
SEC is below 0.1567~kWh/kg, is neutral near that value, and becomes an
increasing loss above it.

\begin{center}\small
\begin{tabularx}{\textwidth}{@{}Y r r r@{}}
\toprule
\textbf{All-in EO SEC} & \textbf{Historical 10\%} & \textbf{O30} &
\textbf{Full duty} \\
& \multicolumn{3}{c}{\textbf{Conservative conditional window [kW]}} \\
\midrule
0.050 kWh/kg & 40.01 & 131.06 & 180.27 \\
0.100 kWh/kg & 21.26 & 69.64 & 95.80 \\
0.150 kWh/kg & 2.51 & 8.23 & 11.32 \\
\bottomrule
\end{tabularx}
\end{center}

The exact endpoint has a calculated 0.2752-kWh/kg break-even because the
remaining coil becomes sensible-only. Infinitesimally above the endpoint, the
screening logic returns to the latent-control branch and the conservative
budget is 0.1567~kWh/kg. Use the lower one-sided budget for decisions until a
real coil/control sequence validates the ideal endpoint transition.

\reportfigure[0.95]{../../outputs/post_erv_eo_percentage_sweep_v1/percentage_membrane_receiver.png}{
Fixed-coordinate membrane area, voltage and receiver activity versus
remaining-duty completion. This is a re-sized trajectory at each target, not
one fixed installed module.}
\FloatBarrier

\subsection{Full-duty uniform-current sensitivity}

The best target-reaching balanced-receiver row in the uniform-current study
uses 8,376.2~m\textsuperscript{2}, 3.0~A/m\textsuperscript{2} and
1.895~V. It leaves 216.80~kW on the conservative near-endpoint branch before
receiver/sink work, or 417.00~kW on the ideal exact-endpoint branch. The
receiver/sink SEC room is 0.1283~kWh/kg. This is useful as a thermodynamic and
sizing sensitivity, not evidence that one electrical bus or a practical
receiver can realise the point.

\reportfigure[0.94]{../../outputs/post_erv_full_eo_study_v1/receiver_and_sink_budget.png}{
Conditional receiver/sink SEC budget for full post-ERV duty. Every receiver,
regeneration, disposal, pumping and heat-rejection term must fit inside the
remaining budget for the energy sign to survive.}
\FloatBarrier

\section{Common-voltage findings}
\label{sec:commonbus}
\eyebrow{One bus with local current sharing is stricter than a uniform-current sweep}

The common-voltage topology applies one terminal voltage while allowing local
parallel current sharing. No current-density acceptance ceiling and no area
acceptance cap are used. Feasibility requires endpoint, water, voltage and
thermodynamic closure.

\begin{center}\small
\begin{tabularx}{\textwidth}{@{}Y r r r r Y@{}}
\toprule
\textbf{Replay} & \textbf{Completion} & \textbf{Voltage} & \textbf{Area} &
\textbf{$j$ range} & \textbf{Status} \\
& [\%] & [V] & [m\textsuperscript{2}] & [A/m\textsuperscript{2}] & \\
\midrule
50\% & 50.000 & 1.000 & 2,468.35 & 3.646--6.879 & Converged target \\
60\% & 60.000 & 1.000 & 3,708.38 & 2.432--6.322 & Converged target \\
67.5\% & 67.500 & 1.000 & 5,269.67 & 1.520--5.905 & Converged target \\
O30 & 72.700 & 1.000 & 7,191.41 & 0.887--5.616 & Converged target \\
O30 balanced-anchor voltage & 72.700 & 1.314 & 3,461.40 & 3.044--8.427 & Converged target \\
Full duty & 100.000 & 1.000 & -- & 0.064--10.964 & No physical common-voltage solution at this coordinate \\
Full duty & 100.000 & 2.000 & -- & 2.358--10.006 & Numerical/local failure; not a no-go proof \\
\bottomrule
\end{tabularx}
\end{center}

\reportfigure[0.96]{../../outputs/project_closure_check_v1/figures/closure_common_voltage.pdf}{
Project-closure replay of the common-voltage window. O30 and lower completion
targets have confirmed sampled points. Full duty has no confirmed sampled
point in the frozen range and remains an open architecture question.}

\subsection{Diffusivity sensitivity}

The refined one-factor study varies $D$ while freezing the remaining
inverse-design properties. O30 has confirmed sampled voltages at 1.00--1.25~V
for all seven $D$ coordinates. At $D=10^{-14}$~m\textsuperscript{2}/s, the
selected 1.00-V row uses approximately 7,174~m\textsuperscript{2}, consumes
30.06~kW known EO+fan power and leaves a conditional 162.42-kW window. At
$D=10^{-9}$~m\textsuperscript{2}/s, the selected known power rises to
373.89~kW and the conditional window becomes $-181.41$~kW.

\reportfigure[0.96]{../../outputs/workshop_common_voltage_diffusivity_v1/figures/workshop_common_voltage_diffusivity.pdf}{
Common-voltage diffusivity cross-check. Crosses denote sampled physical
no-solutions and open triangles unresolved numerical/local high-bracket rows.
No sampled value is not a global proof of impossibility.}

For full remaining duty, none of 49 refined sampled voltages is confirmed:
21 are physical no-solutions and 28 are numerical/local unresolved. Across
both segment counts, 98 flagged rows were retried at a 100-times higher
configured numerical bound; none recovered full closure. The allowed wording
is therefore: \emph{no confirmed feasible sampled common-voltage point was
found below 2~V at the frozen coordinate}. Voltage zoning, staging or changed
receiver/path physics belongs after measured material data.

\section{Figure of merit and energy budget}
\eyebrow{FOM is the reciprocal of SEC only after the power boundary is named}

\[
\fom=\frac{\dot m_{\mathrm{water}}}{P_{\mathrm{device}}}
\quad[\mathrm{kg/kWh}],
\qquad
\secpd=\frac{P_{\mathrm{device}}}{\dot m_{\mathrm{water}}}
\quad[\mathrm{kWh/kg}],
\qquad
\fom=\secpd^{-1}.
\]

\begin{center}\small
\begin{tabularx}{\textwidth}{@{}Y r r r r Y@{}}
\toprule
\textbf{Device coordinate} & \textbf{Water} & \textbf{Known power} &
\textbf{Known SEC} & \textbf{Known FOM} & \textbf{Boundary} \\
& [kg/h] & [kW] & [kWh/kg] & [kg/kWh] & \\
\midrule
Post-ERV EO to O30, 1.00 V & 1,228.27 & 30.057 & 0.02447 & 40.864 &
EO grid plus idealised module fan \\
Direct EO in ERV, full duty, 0.25 V & 4,827.15 & 63.419 & 0.01314 & 76.116 &
Reclassified active core plus idealised fan \\
\bottomrule
\end{tabularx}
\end{center}

\reportfigure[0.94]{../../outputs/fom_assessment_v1/figures/device_fom.pdf}{
Known-core FOM and the minimum all-in FOM needed to retain a positive
conditional margin. The open gap between known and break-even performance is
an allowance for missing loads, not evidence that those loads are small.}

For the post-ERV O30 common-voltage point, the minimum break-even all-in FOM is
6.382~kg/kWh and the unallocated SEC budget is 0.13223~kWh/kg. The current
known-core FOM is attractive only if the receiver, thermal and balance-of-plant
terms fit inside that allowance.

Using the complete outdoor-to-supply water duty as a common numerator, the
system water-energy FOMs are 4.139~kg/kWh for AC only, 8.878 for passive
ERV+AC, 12.659 for post-ERV EO O30, and 44.755 for the direct active-ERV
full-duty screen. The EO system values remain conditional incomplete ledgers.
Twelve of thirteen requested FOMs can be calculated or parameterised; payback
cannot be stated without CAPEX and incremental maintenance.

\section{Membrane and material findings}
\eyebrow{The target is a coupled property vector, not one favourable scalar}

\subsection{Direct EO and post-ERV parameter regions}

The direct-EO 2-V screen has 178 membrane-gate passes among 2,508 physical
coordinates; 108 remain at the preferred 1.8-V boundary. Those are discrete
sampled supports, not continuous tolerance boxes. With a 0.85-V combined fixed
loss, none is positive against the modern comparator before receiver energy.
A deterministic 0.20-V fixed-loss scenario has 15 positive rows, still before
receiver and balance-of-plant terms.

For EO after the ERV, 646 of 1,122 dense $D$--$j$ rows are conditionally
energy-positive on the declared coefficient basis. This does not imply that
646 materials exist: all other properties were frozen.

\reportfigure[0.90]{../../outputs/material_parameter_windows_with_nafion_v1/04_simultaneous_property_gap_vs_nafion.pdf}{
Simultaneous property gap relative to the declared Nafion reference. Marginal
axis overlap must not be combined into a fictitious joint material vector.}

\subsection{Nafion interpretation}

Nafion's declared drag envelope is approximately 0.3--2.5 mol
H$_2$O/mol charge, far below the direct-EO drag region in the current screen.
Its literature diffusivity can lie on a sampled post-ERV $D$ axis, but its
conductivity magnitude and activity exponents differ from the frozen
inverse-design model. The adapter has non-negligible fit errors and some
active-ERV local states lie outside its declared hydration domain.

A matched zero-current control is especially instructive. At the selected
Nafion-reference coordinate, the active row attributes 91.63\% of net water
transfer to diffusion and consumes 0.444~kW for EO. The exact passive control
requires only 2.653\% more area and uses 0.493~kW less electricity.
Consequently, that apparent benefit is primarily passive membrane transfer,
not an EO energy result.

\begin{warningbox}
\textbf{Material non-gates.} Nafion is not the target. There is no current
density acceptance gate at 21~A/m\textsuperscript{2}, area is an observation
rather than a pass/fail cap in this phase, and a drag number near 100 is not
sufficient by itself. Do not mix properties from different samples,
conditioning states or literature schemas.
\end{warningbox}

\subsection{Same-sample data required before ranking}

\begin{center}\small
\begin{tabularx}{\textwidth}{@{}p{19mm}p{48mm}Y@{}}
\toprule
\textbf{ID} & \textbf{Property or measurement} & \textbf{Decision role} \\
\midrule
MAT-01 & Sorption isotherm, hysteresis and water capacity & Converts RH into
chemical loading/work and defines the valid domain \\
MAT-02 & $D(a,T)$ or chemical mobility/permeance with definition & Separates
kinetics from capacity and quantifies passive co-/back-diffusion \\
MAT-03 & EO drag $n_d(a,T,j)$ & Sets polarity-odd water transported per charge \\
MAT-04 & Wet through-plane conductivity & Sets bulk ohmic loss \\
MAT-05 & Compressed membrane plus contact ASR & Prevents optimism from bulk
conductivity alone \\
MAT-06 & Dry/wet thickness, swelling, creep and mechanics & Couples permeance,
resistance, gap control and lifetime \\
MAT-07 & Coupled two-sided net water flux & Tests whether the property bundle
coexists under model-relevant gradients \\
MAT-08 & Receiver/electrode wetting, corrosion, flooding and drainage & Closes
interface compatibility and regeneration needs \\
MAT-09 & Sample, batch, ion form, conditioning and raw-data provenance & Makes
the candidate auditable and procurable \\
\bottomrule
\end{tabularx}
\end{center}

Post-ERV polishing and direct-ERV integration can prefer different materials.
Post-ERV EO benefits from suppressing adverse chemical permeance. Direct ERV
can exploit favourable passive transfer before a pinch, so minimising
permeance indiscriminately can remove useful free transport.

\section{EO vapour coupon-cell dimensions}
\eyebrow{Centimetre-scale, controlled-flow and reconfigurable}

\begin{center}\small
\begin{tabularx}{\textwidth}{@{}p{47mm}p{43mm}Y@{}}
\toprule
\textbf{Element} & \textbf{Starting value} & \textbf{Reason} \\
\midrule
Maximum active aperture & 50$\times$50 mm = 25 cm\textsuperscript{2} &
Maximises signal while retaining a manageable fixture \\
Additional square masks & 30$\times$30 and 20$\times$20 mm &
Area scaling and lower-signal checks \\
Identifiability inserts & 40$\times$40 mm and equal-area 8$\times$50 mm &
Separates area/perimeter and orientation effects \\
Membrane coupon & Approximately 80$\times$80 mm &
15-mm land per side for seals, guard and handling \\
Outer fixture & Approximately 110$\times$110 mm &
Additional land for bolts and ports \\
Symmetric gaps & 1, 2 and 4 mm; begin at 2 mm &
Direct film and boundary sensitivity \\
Flow frame & 8--10 mm insulating frame &
Starting structural and port depth \\
Backing/load plate & 12 mm metal; 16 mm if load-bearing PEEK &
Conservative deflection screen pending measured torque/flatness \\
Nominal high-flow mode & 6 actual L/min per side at 2-mm gap &
Approximately 1 m/s, Re about 236, ideal straight-channel drop about 2.9 Pa \\
Swept volume & 5.00 mL active; 6.22 mL with headers/ports &
Keep added dead volume at or below about 1.3 mL per side \\
\bottomrule
\end{tabularx}
\end{center}

\reportfigure[0.96]{../../outputs/eo_vapor_cell_sizing_v1/cell_dimension_tradeoffs.png}{
Coupon-cell dimension trade-offs. The 25-cm\textsuperscript{2}, 2-mm-gap
fixture is the nominal starting point; interchangeable masks and gaps are part
of the measurement design, not optional refinements.}

The plate screen limits bow to 20~$\mu$m, or 2\% of the 1-mm gap, across a
50-mm unsupported square. Gas pressure alone is not governing. Under a
deliberately conservative 1-MPa equivalent uniform fixture-load surrogate,
calculated minimum thicknesses are 15.4~mm PEEK, 5.8~mm aluminium and
4.2~mm 316L. These are screening values, not a gasket-contact calculation;
the assembled cell needs a torque, loaded-gap and flatness map.

\subsection{Signal and flow modes}

The 25-cm\textsuperscript{2} gross EO signal screen spans approximately
5.04--840.2~mg/h. With 0.1 RH-percentage-point expanded uncertainty for each
humidity channel and a signal-to-uncertainty gate of five, low-drag signals
would require flows below the 6-L/min film-characterisation condition. Use
separate modes:
\begin{itemize}
  \item \textbf{high-flow film mode:} nominally 6~L/min per side, repeated at
        3 and 12~L/min and at 1/2/4-mm gaps;
  \item \textbf{low-flow/integrating signal mode:} approximately
        0.69~L/min for the 5.04-mg/h floor or 2.29~L/min for the
        16.80-mg/h order-unity reference, unless a more sensitive integrating
        water measurement is used.
\end{itemize}

A shallow closed chamber cannot hold fixed RH under transport. In the
gas-only 5-mL chamber, the low-drag 5.04-mg/h reference reaches a 5-Pa
differential in 0.063~s. Even a 5-L reservoir extends that ideal time only to
63.4~s. Closed operation is therefore a transient inventory experiment with
pressure compliance/control, not a steady boundary.

\section{Boundary sensitivity and error control}
\eyebrow{Polarity reversal removes static bias, not films or dynamic drift}

\reportfigure[0.95]{../../outputs/eo_vapor_boundary_analysis_v1/combined_sensitivity_rank.png}{
Ranked boundary and nuisance sensitivities in the declared stress space. The
ranges are experiment-design scenarios, not calibrated probability
distributions or confidence limits.}

At the nominal 25-cm\textsuperscript{2}, 2-mm, 32$^\circ$C condition:
\begin{itemize}
  \item a static +0.1 RH-percentage-point inlet mismatch changes the
        fast-diffusion one-way result by about +10.2\%, while the ideal
        polarity-odd estimator nearly rejects it;
  \item a +0.1-K side-to-side temperature mismatch changes the same one-way
        result by about $-28.9$\%, while a stable reversal again nearly
        rejects the static part;
  \item changing the gap from 2 to 1~mm shifts the fast-diffusion odd signal
        by about $-2.0$\%; changing to 4~mm shifts it by about $-10.1$\%;
  \item high-diffusion cases are dominated by the unresolved film multiplier,
        while slow-diffusion cases are dominated by humidity drift between
        polarities.
\end{itemize}

ABBA ordering rejects static and linear drift only when the boundary is stable
over the block. Polarity-correlated humidity, wall inventory, nonlinear drift,
hysteresis and settling can survive the odd estimator and mimic EO. Therefore:
\begin{enumerate}
  \item feed both sides from a split common inlet for equal-RH work;
  \item cross-switch the humidity reference and randomise or mirror ABBA block
        order;
  \item record source loss and receiver gain independently and require water
        closure within combined uncertainty;
  \item control membrane differential pressure to the declared $\pm5$-Pa
        target and deliberately challenge it separately;
  \item run impermeable-foil blanks before and after each block;
  \item report one-way, polarity-even and polarity-odd rates separately;
  \item repeat flow and gap, and fit area plus perimeter using at least five
        inserts including an equal-area/different-perimeter pair.
\end{enumerate}

The original three square masks alone give zero residual degrees of freedom
for an area-plus-perimeter fit and an area/perimeter correlation of 0.995.
The recommended five-insert set yields two residual degrees of freedom and
reduces, but does not eliminate, the correlation.

\section{Open gates}
\eyebrow{What still prevents a design or energy claim}

\begin{center}\small
\begin{tabularx}{\textwidth}{@{}p{40mm}p{22mm}Y@{}}
\toprule
\textbf{Gate} & \textbf{State} & \textbf{Required closure evidence} \\
\midrule
Candidate data completeness & \textcolor{Orange}{\textbf{OPEN}} &
One same-sample MAT-01--MAT-09 dataset with raw points and provenance \\
Receiver boundary and SEC & \textcolor{Orange}{\textbf{OPEN}} &
Chosen gas/liquid/desiccant sink; measured transfer, pumping, regeneration or
disposal and compatibility \\
EO heat rejection & \textcolor{Orange}{\textbf{OPEN}} &
Conjugate energy balance with transported-water enthalpy assigned once \\
Full-duty common bus & \textcolor{Orange}{\textbf{OPEN}} &
Measured candidate replay and, if warranted, staged/zoned topology \\
Films and interfaces & \textcolor{Orange}{\textbf{OPEN}} &
Flow/gap/temperature-resolved source and receiver exchange coefficients \\
Hardware/BOP & \textcolor{Orange}{\textbf{OPEN}} &
Electrodes, contacts, converter topology, manifolds, sealing, leakage and controls \\
Practical module limits & \textcolor{Orange}{\textbf{OPEN}} &
Team-defined area, footprint, pressure drop, CAPEX, maintenance and lifetime gates \\
Durability and safety & \textcolor{Orange}{\textbf{OPEN}} &
Wet cycling, creep, corrosion, fouling, electrical safety and long-duration mass balance \\
Immutable release & \textcolor{Orange}{\textbf{OPEN}} &
Freeze reviewed source/output hashes in a clean commit or equivalent snapshot \\
Technology validation & \textcolor{Red}{\textbf{NOT ACHIEVED}} &
Measured coupled material, subsystem and full-system evidence \\
\bottomrule
\end{tabularx}
\end{center}

\section{Recommended next experiments}
\eyebrow{Convert the remaining uncertainty into measurements}

\subsection{Experiment 1: fixture qualification and blank closure}

\begin{enumerate}
  \item Assemble the 50$\times$50-mm aperture cell with 2-mm gaps, metal load
        plates and double seals.
  \item Map loaded gap and flatness against bolt torque; measure swept and dead
        volumes.
  \item Run foil blanks over 3/6/12~L/min, 1/2/4-mm gaps and the intended
        pressure envelope.
  \item Accept the fixture only when leakage, pressure-driven bypass and wall
        inventory fit within a predeclared fraction of the minimum useful
        polarity-odd signal.
\end{enumerate}

\subsection{Experiment 2: EO identification with matched zero-current control}

\begin{enumerate}
  \item Use a split common inlet and stabilise temperature, RH and pressure.
  \item Run $+J/0/-J$ and mirrored ABBA blocks at multiple current levels.
  \item Measure water on both sides, charge, terminal voltage and local
        temperatures; report odd, even and one-way terms.
  \item Repeat with flow, gap, aperture area and perimeter changes.
  \item Treat any current-off transport as the matched passive control, not as
        background to be assumed negligible.
\end{enumerate}

\subsection{Experiment 3: same-sample property bundle}

For every candidate, freeze sample ID, batch, ion form, conditioning and wet
thickness. Measure sorption/hysteresis, activity-resolved chemical mobility,
EO drag, wet conductivity, compressed ASR, swelling/mechanics and the coupled
two-sided flux. Candidate ranking begins only when all MAT-01--MAT-09 rows are
complete for the same sample.

\subsection{Experiment 4: receiver and thermal ledger}

Choose one receiver concept and measure or bound water acceptance, pressure
drop, pumping, regeneration/disposal, flooding/drainage, electrode
compatibility and heat rejection. The pass criterion is not receiver SEC
alone: the complete added ledger must remain inside the O30 conditional
0.13223-kWh/kg allowance with an uncertainty margin.

\subsection{Experiment 5: measured-data topology replay}

Only after Experiments 1--4, replay:
\begin{itemize}
  \item post-ERV O30 on one common bus;
  \item higher completion with one bus;
  \item two voltage zones or staged EO if full-duty conflict remains;
  \item direct active EO in the ERV if measured passive transfer and films
        justify it.
\end{itemize}
The replay must include measured property covariance and uncertainty, not
independent best-case axes.

\begin{mainbox}
\textbf{Recommended sequence.} Fixture qualification $\rightarrow$ matched
passive/active coupon identification $\rightarrow$ same-sample property bundle
$\rightarrow$ receiver/thermal ledger $\rightarrow$ measured-data topology
replay. This sequence closes a decision gate at every step and avoids another
unfocused model sweep.
\end{mainbox}

\section{Provenance and claim boundary}
\eyebrow{Canonical machine-readable outputs, not copied report assets}

This v8 report references the current figures and machine-readable evidence
in place; no output figure or data file is copied into this report directory.
The principal sources are:

\begin{itemize}
  \item \evidencepath{outputs/project_closure_check_v1/closure_status.json},
        \evidencepath{outputs/project_closure_check_v1/decision_gate_register.csv}
        and \evidencepath{outputs/project_closure_check_v1/source_manifest.json};
  \item \evidencepath{outputs/final_delivery_check_v1/status.json},
        \evidencepath{outputs/final_delivery_check_v1/gate_register.csv} and
        \evidencepath{outputs/final_delivery_check_v1/source_manifest.json};
  \item \evidencepath{outputs/post_erv_full_eo_study_v1/}:
        \evidencepath{SUMMARY.md}, energy windows, transport rows and audit;
  \item \evidencepath{outputs/post_erv_eo_percentage_sweep_v1/}:
        percentage semantics, energy ledger, fixed-coordinate trajectory and audit;
  \item \evidencepath{outputs/workshop_common_voltage_diffusivity_v1/}:
        case register, refined envelope, bracket probe, profiles, audit and source manifest;
  \item \evidencepath{outputs/fom_assessment_v1/system_fom.csv},
        \evidencepath{outputs/fom_assessment_v1/device_fom.csv},
        \evidencepath{outputs/fom_assessment_v1/summary.json},
        \evidencepath{outputs/fom_assessment_v1/audit.json} and
        \evidencepath{outputs/fom_assessment_v1/source_manifest.json};
  \item \evidencepath{outputs/material_parameter_windows_with_nafion_v1/} and
        \evidencepath{outputs/direct_passive_nafion_erv_control_v1/};
  \item \evidencepath{outputs/integrated_evidence_v1/material_data_requirements.csv}
        and \evidencepath{outputs/integrated_evidence_v1/next_study_matrix.csv};
  \item \evidencepath{outputs/eo_vapor_cell_sizing_v1/audit.json},
        \evidencepath{outputs/eo_vapor_cell_sizing_v1/source_manifest.json},
        \evidencepath{outputs/eo_vapor_boundary_analysis_v1/audit.json} and
        \evidencepath{outputs/eo_vapor_boundary_analysis_v1/source_manifest.json}.
\end{itemize}

\subsection{Permitted conclusions}

\begin{itemize}
  \item The model identifies conditional post-ERV and direct-ERV energy windows.
  \item O30 is a confirmed sampled common-voltage mechanism coordinate and the
        preferred first system-relevant coupon target.
  \item O50 needs no post-ERV EO duty at the declared absolute-humidity basis.
  \item Full post-ERV duty on one common bus below 2~V is not demonstrated in
        the sampled frozen coordinate and is not proven impossible.
  \item The material question is whether a coupled property bundle exists in
        one conditioned and durable sample.
  \item The model phase is sufficiently closed to begin material and subsystem
        feasibility work.
\end{itemize}

\subsection{Prohibited conclusions}

\begin{itemize}
  \item No listed kW difference is a realised or validated saving.
  \item No FOM is an all-in product FOM unless every omitted term is added.
  \item No sampled material coordinate proves a continuous tolerance or a
        commercially available membrane.
  \item No current-density or area value is an acceptance limit in the present
        exploratory phase.
  \item Audit PASS is not material, module, system or commercial validation.
\end{itemize}

\begin{stopbox}
The report is a current decision synthesis, not an immutable release. Before
external distribution, freeze the reviewed repository state, rerun the final
delivery check and all declared tests, regenerate the PDF, record hashes and
package the source with an explicit manifest.
\end{stopbox}

\appendix
\section{Compact numerical register}

\begin{longtable}{@{}p{47mm}p{29mm}p{27mm}p{57mm}@{}}
\toprule
\textbf{Metric} & \textbf{Value} & \textbf{Evidence class} & \textbf{Canonical source} \\
\midrule
\endfirsthead
\toprule
\textbf{Metric} & \textbf{Value} & \textbf{Evidence class} & \textbf{Canonical source} \\
\midrule
\endhead
AC-only equal-service power & 1,166.400 kW & reference anchor &
\evidencepath{outputs/fom_assessment_v1/system_fom.csv} \\
Passive ERV+AC power & 543.740564 kW & modelled reference &
\evidencepath{outputs/fom_assessment_v1/system_fom.csv} \\
Post-ERV O30 conditional subtotal & 381.326719 kW & incomplete ledger &
\evidencepath{outputs/fom_assessment_v1/system_fom.csv} \\
Post-ERV O30 conditional window & 162.413845 kW & incomplete ledger &
\evidencepath{outputs/fom_assessment_v1/system_fom.csv} \\
Post-ERV O30 known-core FOM & 40.8641 kg/kWh & known core only &
\evidencepath{outputs/fom_assessment_v1/device_fom.csv} \\
O30 common-voltage area at FOM point & 7,173.55 m\textsuperscript{2} & model coordinate &
\evidencepath{outputs/fom_assessment_v1/device_fom.csv} \\
O30 all-in break-even SEC & 0.156701 kWh/kg & conditional budget &
\evidencepath{outputs/post_erv_full_eo_study_v1/break_even_by_target.csv} \\
O30 remaining-duty completion & 72.7003\% & exact coordinate &
\evidencepath{outputs/post_erv_eo_percentage_sweep_v1/percentage_coordinates.csv} \\
Historical 10\% remaining-duty completion & 22.1925\% & traceability coordinate &
\evidencepath{outputs/post_erv_eo_percentage_sweep_v1/percentage_coordinates.csv} \\
Full-duty conservative sink budget & 0.128323 kWh/kg & conditional budget &
\evidencepath{outputs/project_closure_check_v1/closure_status.json} \\
O30 common-voltage diffusivity span & $10^{-14}$--$10^{-9}$ m\textsuperscript{2}/s sampled &
one-factor screen &
\evidencepath{outputs/workshop_common_voltage_diffusivity_v1/diffusivity_envelope.csv} \\
Full-duty refined sampled outcomes & 21 physical / 28 unresolved & open topology gate &
\evidencepath{outputs/workshop_common_voltage_diffusivity_v1/diffusivity_envelope.csv} \\
Active coupon aperture & 50$\times$50 mm & starting fixture &
\evidencepath{outputs/eo_vapor_cell_sizing_v1/fixture_dimension_basis.csv} \\
Nominal coupon gap and flow & 2 mm / 6 L min\textsuperscript{-1} side\textsuperscript{-1} &
film mode &
\evidencepath{outputs/eo_vapor_cell_sizing_v1/channel_geometry_envelope.csv} \\
\bottomrule
\end{longtable}

\medskip
{\sans\scriptsize\color{Slate}
End of report. Current interpretation as of 31 August 2026. All EO system
comparisons remain conditional until measured material, receiver, thermal and
hardware ledgers close.}

\end{document}
